\documentclass[10pt,a4paper]{amsart}
\usepackage[margin=19mm]{geometry}
\usepackage{amsmath,amssymb,mathtools}
\usepackage[scr=rsfs,cal=euler]{mathalfa}
\usepackage{xfrac}
\usepackage[unicode,hidelinks,bookmarks=false]{hyperref}
\newcommand{\ie}{i.e.\ }
\newcommand{\cf}{cf.\ }
\newcommand{\resp}{respectively\ }
\newcommand{\Subsection}{\subsection}
\providecommand{\nobreakdash}{\nobreak\hbox{-}}
\setlength{\emergencystretch}{2em}
\multlinegap0pt
\usepackage{amssymb}
\usepackage{mathtools}
\usepackage{bm}
\usepackage[shortlabels]{enumitem}
\newtheorem{lemma}{Lemma}
\usepackage{cleveref}
\crefname{lemma}{Lemma}{Lemmas}
\title{Improved Analysis of the Accelerated Noisy Power Method with Applications to Decentralized PCA}
\author{Pierre Agui\'e, Mathieu Even, Laurent Massouli\'e}
\date{Source: arXiv:2602.03682v2 (CC BY 4.0)}
\begin{document}
\maketitle

\noindent\emph{Source and licence.} Improved Analysis of the Accelerated Noisy Power Method with Applications to Decentralized PCA. Version arXiv:2602.03682v2 (CC BY 4.0), distributed by arXiv under the Creative Commons Attribution 4.0 International licence, http://creativecommons.org/licenses/by/4.0/, as recorded in the arXiv licence metadata for this exact version and shown on the abstract page https://arxiv.org/abs/2602.03682v2. Full licence notice: \texttt{licenses/arxiv-2602.03682-CC-BY-4.0.txt}.\par\smallskip

\noindent\textit{Editorial scope.} Counted unit: the article's lemma lemma:ht\_exp\_decay (source0.tex lines 1286--1291). The article's complete original proof of it is delivered with it (source0.tex lines 1293--1387, one \texttt{proof} environment of 95 lines). No mathematics is written, repaired or re-derived: the statement and the proof are the article's own lines, in the article's own notation, and no retained line is substituted or re-flowed.  Retained uncounted as the source-local context the proof needs: lines 1247--1249; lines 1406--1412.  Nothing was omitted from any retained span, so the package carries no omission record.  No retained line contains a citation, so the wrapper carries no bibliography and no bibliography entry.  No retained line contains a figure environment, an image-inclusion command or drawing code, so no image is delivered and the package declares no asset dependency.  Editorial formatting only, in addition: an AMS \texttt{amsart} preamble that restates the article's own declarations and macros used by the retained lines, namely the environments \texttt{lemma} on separate counters, together with the standard packages those lines need.  The retained environments print in this document as Lemma 1 in order of appearance, whereas the article prints the same items with the numbering of its own full document.  The complete native source of the article as distributed in the arXiv e-print for this exact version is bundled unchanged as \texttt{sources/193924/source0.tex}.
  \begin{align}\label{eq:ht_reccurence}
    h_t & \leqslant c_1\gamma^t h_0 + \eta \sum_{s=0}^{t-1} (s+1) \gamma^s (1 + h_{t-1-s}),
  \end{align}

\begin{lemma}\label{lemma:f_delta_bounds}
  For all $\Delta \in [0,1)$, let
  \begin{align*}
    f(\Delta) := \frac{\sqrt{2-\Delta} - \sqrt{\Delta}}{1 - \Delta}.
  \end{align*}
  Then, f is decreasing over $[0,1)$. In particular, we have $1 \leqslant f(\Delta) \leqslant \sqrt{2}$ for all $\Delta \in [0,1)$, and $f(\Delta) \leqslant \sqrt{3} - \sqrt{2} < 1$ for all $\Delta \in [1/2, 1)$.
\end{lemma}

\begin{lemma}\label{lemma:ht_exp_decay}
    For all $t\geqslant 0$,
  \begin{align}\label{eq:ht_exp_decay}
    h_t \leqslant \varepsilon / 2 + h_0 \left(1 - \sqrt{\Delta}/2\right)^t.
  \end{align}
\end{lemma}

\begin{proof}
  Let $\{y_t\}$ be the sequence defined as $y_0 = h_0$ and for all $t\geqslant 1$,
  \begin{align}\label{eq:yt_def}
    y_t & = c_1 \gamma^t h_0 + \eta \sum_{s=0}^{t-1} (s+1) \gamma^s (1 + y_{t-1-s}).
  \end{align}
  Then, for all $t \geqslant 0$, we have $h_t \leqslant y_t$ (this can be shown by a simple induction using the recurrence inequality \eqref{eq:ht_reccurence}). The proof will then consist in upper bounding $y_t$ by the right-hand side of \eqref{eq:ht_exp_decay}.

  To do so, consider the generating function of the sequence $\{y_t\}$ defined as $Y(z) := \sum_{t=0}^{+\infty} y_t z^t$. We will use the following identities, which are true for all $\vert z \vert < 1/\gamma$:
  \begin{align*}
    \sum_{t=0}^{+\infty} \gamma^t z^t & = \frac{1}{1 - \gamma z}, \quad \quad
    \sum_{t=0}^{+\infty} (t+1) \gamma^t z^t = \frac{1}{(1 - \gamma z)^2}.
  \end{align*}
  
  Then, by multiplying both sides of \eqref{eq:yt_def} by $z^t$ and summing over all $t\geqslant 1$, we have
  \begin{align*}
    Y(z) - y_0 & = c_1 h_0 \sum_{t=1}^{+\infty} \gamma^t z^t + \eta z \sum_{t=1}^{+\infty} \sum_{s=0}^{t-1} (s+1) \gamma^s z^{t-1} + \eta z \sum_{t=1}^{+\infty} \sum_{s=0}^{t-1} (s+1) \gamma^s y_{t-1-s} z^{t-1}, \\
    Y(z) - h_0 & = c_1 h_0 \left( \frac{1}{1 - \gamma z} - 1 \right) + \eta z\frac{1}{(1 - \gamma z)^2} \frac{1}{1-z} + \eta z Y(z) \frac{1}{(1 - \gamma z)^2}, \\
    Y(z) \left( 1 - \frac{\eta z}{(1 - \gamma z)^2} \right) & = h_0 \left( 1 + c_1  \frac{\gamma z}{1 - \gamma z}  \right) + \eta z \frac{1}{(1 - \gamma z)^2}\frac{1}{1-z}, \\
    Y(z) & = \frac{h_0 \left( 1 + c_1 \frac{\gamma z}{1 - \gamma z} \right) + \eta z \frac{1}{(1 - \gamma z)^2}\frac{1}{1-z}}{1 - \frac{\eta z}{(1 - \gamma z)^2}}.
  \end{align*}
  Here, the second inequality used the fact that the generating function of the convolution of two sequences is the product of the two associated generating functions. These equalities are true for all $z$ whose magnitude is less than the pole of the right-hand side with lowest magnitude, which is non-zero. Indeed, the poles of $Y(z)$ are $1$ and the roots of the polynomial $(1-\gamma z)^2 - \eta z$, which are
  \begin{align*}
    \rho_\pm := \frac{2\gamma + \eta \pm \sqrt{(2\gamma + \eta)^2 - 4\gamma^2}}{2\gamma^2} > 0 .
  \end{align*}

  We can then decompose $Y(z)$ into partial fractions:
  \begin{align}\label{eq:yz_partial}
    Y(z) = \frac{\kappa_1}{1-z} + \frac{\kappa_+}{1 - z/\rho_+} + \frac{\kappa_-}{1 - z/\rho_-},
  \end{align}
  from which we deduce that for all $t\geqslant 0$,
  \begin{align}
    y_t & = \kappa_1 + \kappa_- \left(\frac{1}{\rho_-}\right)^t + \kappa_+ \left(\frac{1}{\rho_+}\right)^t \leqslant \kappa_1 + (\kappa_+ + \kappa_-)\left(\frac{1}{\rho_-}\right)^t\nonumber\\
    h_t &  \leqslant \kappa_1 + (\kappa_+ + \kappa_-)\left(\frac{1}{\rho_-}\right)^t \label{eq:ht_upper}.
  \end{align}

  Here, the first inequality is due to the fact that $\kappa_+$ is non-negative. Indeed, multiplying both sides of \eqref{eq:yz_partial} by $(1 - z/\rho_+)$ and taking the limit $z \to \rho_+$ gives after simplification
  \begin{align*}
    \kappa_+ & = \frac{h_0(1-\gamma\rho_+)(1-\gamma\rho_+ + c_1\gamma\rho_+) + \eta\rho_+/(1-\rho_+)}{1 - \rho_+/\rho_-},
  \end{align*}
  which is non-negative since $\rho_+ > \rho_-$, $1 - \gamma \rho_+ < 0$, $1 - \rho_+ < 0$ and $c_1 \geqslant 1$.

  Multiplying both sides of the partial fraction decomposition of $Y(z)$ in \eqref{eq:yz_partial} by $(1-z)$ and taking the limit $z \to 1$ gives
  \begin{align}\label{eq:kappa1def}
    \kappa_1 & = \frac{\eta}{(1-\gamma)^2 - \eta}.
  \end{align}

  Evaluating \eqref{eq:yz_partial} at $z=0$ gives
  \begin{align}\label{eq:kappa_plusminus}
    \kappa_+ + \kappa_- & = h_0 - \kappa_1.
  \end{align}

  We now analyze in detail the term $1-\gamma$. We have
  \begin{align}
    1-\gamma & = 1 - \frac{\sqrt{\beta}}{(1-c)\lambda_k^+ + c\sqrt{\beta}} \nonumber\\
    & \geqslant 1 - (1-c)\frac{\sqrt{\beta}}{\lambda_k^+} - c \nonumber\\
    & = (1-c)\left(1 - \frac{\sqrt{\beta}}{\lambda_k^+}\right), \label{eq:1-gamma_sqrtbeta}
  \end{align}
  where the inequality is due to the convexity of $u\mapsto 1/u$. Rewriting $1-\sqrt{\beta}/\lambda_k^+$ in terms of $\Delta = 1 - 2\sqrt{\beta}/\lambda_k$ gives
  \begin{align}\label{eq:1_gamma_fdelta}
    1 - \frac{\sqrt{\beta}}{\lambda_k^+} & = \frac{\lambda_k + \sqrt{\lambda_k^2 - 4\beta} - 2\sqrt{\beta}}{\lambda_k + \sqrt{\lambda_k^2 - 4\beta}} = \sqrt{\Delta} \underbrace{\frac{\sqrt{2-\Delta} - \sqrt{\Delta}}{1 - \Delta}}_{=: f(\Delta)}.
  \end{align}
  Then, according to \cref{lemma:f_delta_bounds}, for all $\Delta' \in (0,1)$, we have $f(\Delta' ) \geqslant 1$, so that
  \begin{align*}
    1 - \frac{\sqrt{\beta}}{\lambda_k^+} \geqslant \sqrt{\Delta}.
  \end{align*}

  Plugging this inequality into the previously found lower bound on $(1-\gamma)$ \eqref{eq:1-gamma_sqrtbeta} gives
  \begin{align}\label{eq:1minusgamma_ineq}
    1 - \gamma \geqslant (1-c) \sqrt{\Delta}.
  \end{align}
  Then, $(1 - \gamma)^2 - \eta \geqslant (1-c)^2 \Delta - c_2 \varepsilon \Delta > 0$, so that $\kappa_1 > 0$. We deduce from \eqref{eq:kappa_plusminus} that $\kappa_+ + \kappa_- \leqslant h_0$. Furthermore, from the expression of $\kappa_1$ in \eqref{eq:kappa1def}, we have
  \begin{align*}
    \kappa_1 & = \frac{c_2 \Delta\varepsilon}{(1-\gamma)^2 - c_2 \Delta\varepsilon} \leqslant \frac{c_2 \Delta\varepsilon}{(1-c)^2 \Delta - c_2 \Delta\varepsilon} \leqslant \frac{c_2 \varepsilon}{(1 - c)^2 - c_2} \leqslant \varepsilon/2.
  \end{align*}
  Plugging these inequalities into the upper bound on $h_t$ \eqref{eq:ht_upper} gives
  \begin{align}\label{eq:ht_intermediate}
    h_t & \leqslant  \varepsilon/2 + h_0 \left(\frac{1}{\rho_-}\right)^t.
  \end{align}

  Finally, we analyze the factor $1/\rho_-$. We have
  \begin{align*}
    \frac{1}{\rho_-} & = \gamma \frac{1}{1 + \frac{\eta}{2\gamma} - \sqrt{\left(1 + \frac{\eta}{2\gamma}\right)^2 - 1}} = \gamma \left(1 + \frac{\eta}{2\gamma} + \sqrt{\left(1 + \frac{\eta}{2\gamma}\right)^2 - 1}\right) \\
    & = \gamma + \frac{\eta}{2} + \sqrt{\left( \gamma + \frac{\eta}{2}\right)^2 - \gamma ^2} = \gamma  + \frac{\eta}{2} + \sqrt{\eta \gamma + \frac{\eta^2}{4}} \\
    & \leqslant \gamma + \frac{\eta}{2} + \sqrt{\eta + \eta^2 / 4},
  \end{align*}
  where the last inequality is because $\gamma \leqslant 1$. Since, $\eta = c_2 \varepsilon \Delta \in [0,1]$, we have $\eta^2 \leqslant \eta \leqslant \sqrt{\eta}$, so that
  \begin{align}
    \frac{1}{\rho_-} & \leqslant \gamma + \frac{1 + \sqrt{5}}{2}\sqrt{\eta} \leqslant 1 - (1 - c)\sqrt{\Delta} + \frac{1 + \sqrt{5}}{2}\sqrt{c_2 \varepsilon \Delta} \nonumber\\
    \frac{1}{\rho_-} & \leqslant  1 - \sqrt{\Delta}/2 \label{eq:rho_minus_ineq},
  \end{align}
  where the second equality follows from the inequality of $1-\gamma$ \eqref{eq:1minusgamma_ineq}. Plugging the above bound on $1/\rho_-$ \eqref{eq:rho_minus_ineq} into the inequality on $h_t$ in \eqref{eq:ht_intermediate} gives the desired result:
  \begin{align*}
    h_t & \leqslant \varepsilon/2 + h_0 \left(1 - \sqrt{\Delta}/2\right)^t.
  \end{align*}
\end{proof}

\end{document}
